\documentclass[11pt]{article}

\usepackage[a4paper,margin=25mm]{geometry}
\usepackage{amsmath,amssymb}
\usepackage{booktabs}
\usepackage{hyperref}
\usepackage{microtype}

\newcommand{\dd}{\,\mathrm{d}}

\title{Half-Angle Integration of a Tabulated Phase Function}
\author{}
\date{}

\begin{document}
\maketitle

\section{Phase-function convention}

Let $p(\mu)$ be an azimuthally symmetric phase function, where
\begin{equation}
    \mu = \cos\theta,
    \qquad -1\leq \mu\leq 1,
    \qquad 0\leq\theta\leq\pi.
\end{equation}
The normalization convention used here is
\begin{equation}
    2\pi\int_{-1}^{1}p(\mu)\dd\mu = 1.
    \label{eq:normalization}
\end{equation}
The asymmetry factor is the first angular moment,
\begin{equation}
    g = 2\pi\int_{-1}^{1}\mu p(\mu)\dd\mu.
    \label{eq:g-direct}
\end{equation}
For a strongly forward-peaked phase function, evaluating
Equation~\eqref{eq:g-direct} directly is undesirable: most of the normalization
and most of $g$ come from a narrow region close to $\mu=1$, and $g$ may be very
close to one.

Using Equation~\eqref{eq:normalization}, it is better to compute the small
deficit $1-g$:
\begin{align}
    1-g
    &=2\pi\int_{-1}^{1}p(\mu)\dd\mu
      -2\pi\int_{-1}^{1}\mu p(\mu)\dd\mu \\
    &=2\pi\int_{-1}^{1}(1-\mu)p(\mu)\dd\mu.
    \label{eq:deficit-mu}
\end{align}
The factor $1-\mu$ suppresses the sharp forward peak and makes the moment less
sensitive to errors extremely close to $\mu=1$.

\section{Half-angle coordinate}

Introduce the half-angle variable
\begin{equation}
    x=\sin^2\left(\frac{\theta}{2}\right)
     =\frac{1-\mu}{2},
    \qquad
    \mu=1-2x,
    \qquad
    \dd\mu=-2\dd x.
\end{equation}
The forward direction $\theta=0$ is mapped to $x=0$, while the backward
direction $\theta=\pi$ is mapped to $x=1$.  Reversing the integration limits
gives
\begin{align}
    2\pi\int_{-1}^{1}p(\mu)\dd\mu
        &=4\pi\int_0^1p(1-2x)\dd x, \\
    1-g
        &=8\pi\int_0^1x\,p(1-2x)\dd x.
\end{align}

To obtain still more resolution near the forward direction, set
\begin{equation}
    x=u^q,
    \qquad
    u=x^{1/q}
      =\left(\frac{1-\mu}{2}\right)^{1/q},
    \qquad
    \dd x=q u^{q-1}\dd u,
    \label{eq:u-transform}
\end{equation}
with, for example, $q=2$ or $q=3$.  A uniform grid in $u$ corresponds to
\begin{equation}
    \mu(u)=1-2u^q.
\end{equation}
For $q>1$, this places many corresponding $\mu$ points close to $\mu=1$.

The transformed normalization and deficit are
\begin{align}
    N
    &=4\pi q\int_0^1
      u^{q-1}p(1-2u^q)\dd u,
      \label{eq:N}\\
    D
    &=8\pi q\int_0^1
      u^{2q-1}p(1-2u^q)\dd u.
      \label{eq:D}
\end{align}
Here $D$ is the unnormalized $1-g$ numerator.  If the tabulated phase function
is not exactly normalized, the normalization-corrected asymmetry factor is
\begin{equation}
    \boxed{g=1-\frac{D}{N}}.
    \label{eq:g-corrected}
\end{equation}
When $N=1$, Equation~\eqref{eq:g-corrected} reduces to $g=1-D$.

\section{Interpolation in the transformed coordinate}

Suppose the original table contains angles $\theta_i$ and values $p_i$.  The
coordinates used for interpolation are
\begin{equation}
    \mu_i=\cos\theta_i,
    \qquad
    u_i=\left(\frac{1-\mu_i}{2}\right)^{1/q}.
\end{equation}
The pairs $(u_i,p_i)$ are sorted in increasing $u_i$.  On the interval
$[u_i,u_{i+1}]$, the code uses the linear interpolant
\begin{equation}
    \widetilde p(u)=m_i u+b_i,
    \label{eq:linear-interpolant}
\end{equation}
where
\begin{equation}
    m_i=\frac{p_{i+1}-p_i}{u_{i+1}-u_i},
    \qquad
    b_i=p_i-m_i u_i.
\end{equation}
Transforming the original points before interpolation is important.  Linear
interpolation in $u$ and linear interpolation in $\mu$ define different
approximations between the tabulated samples.  Interpolation in $u$ allocates
more of the representation to the forward-angle structure.

If the first tabulated angle is slightly greater than zero, then $u_0>0$.
The implementation extends the nearest tabulated value as a constant from
$u=0$ to $u_0$.  The same constant-endpoint convention is used near $u=1$ if
the table does not reach $\theta=\pi$.

\section{Exact integration of the linear interpolant}

The transformed normalization integrand can remain sharply peaked even after
the coordinate change.  Instead of resolving that peak with a dense generic
quadrature rule, one can integrate Equation~\eqref{eq:linear-interpolant}
analytically on every tabulated interval.

For any exponent $a>-1$,
\begin{align}
    I_i(a)
    &=\int_{u_i}^{u_{i+1}}u^a\widetilde p(u)\dd u \\
    &=\int_{u_i}^{u_{i+1}}u^a(m_i u+b_i)\dd u \\
    &=\frac{m_i}{a+2}
      \left(u_{i+1}^{a+2}-u_i^{a+2}\right)
      +\frac{b_i}{a+1}
      \left(u_{i+1}^{a+1}-u_i^{a+1}\right).
    \label{eq:interval-integral}
\end{align}
Consequently, the two required quantities are
\begin{align}
    N
    &=4\pi q\sum_i I_i(q-1),
    \label{eq:N-exact}\\
    D
    &=8\pi q\sum_i I_i(2q-1),
    \label{eq:D-exact}\\
    g
    &=1-\frac{D}{N}.
\end{align}
Equations~\eqref{eq:N-exact} and~\eqref{eq:D-exact} are exact for the chosen
piecewise-linear interpolant.  Their accuracy is therefore controlled by the
quality and spacing of the original phase-function table, not by the number of
additional quadrature points.  The number of points used to draw a smooth plot
can be changed without changing $N$, $D$, or $g$.

For comparison, the implementation also permits a uniform-$u$ trapezoidal rule
and Gauss--Legendre quadrature on $u\in[0,1]$.  These are useful convergence
checks, but the interval-analytic method is preferable when the adopted model
is piecewise linear in $u$.

\section{Quantities shown in the four-panel plot}

The plotting script displays the following quantities:
\begin{center}
\begin{tabular}{ll}
\toprule
Panel & Quantity \\
\midrule
Upper left & $p(\mu)$ and its linear interpolation in $\mu$ \\
Upper right & $p(1-2u^q)$ and its linear interpolation in $u$ \\
Lower left & $4\pi q\,u^{q-1}p(1-2u^q)$, whose integral is $N$ \\
Lower right & $8\pi q\,u^{2q-1}p(1-2u^q)$, whose integral is $D$ \\
\bottomrule
\end{tabular}
\end{center}
The extra power of $u^q$ in the deficit integrand explains why its forward
peak is much more strongly suppressed than the normalization integrand.

\end{document}
